%% toc: P2 21 21.8
\section{Finite versus unrestricted validity and countermodel classes}
\label{sec:p02-w02-p02-036-finite-validity}

\emph{Purpose.}  Unrestricted first-order validity is computably enumerable,
because a sound and complete effective calculus lists the valid sentences
(\cref{thm:fol-validity-enumerable}).  This section shows that nothing of the
kind holds when truth is required only in the finite structures.  It sorts
the sentences of a fixed finite relational signature into three
countermodel classes, according to whether they have a finite model, only
infinite models, or no model at all.  It proves Trakhtenbrot's theorem
(\cref{thm:trakhtenbrot}) in full, by encoding the accepting computations of
a Turing machine as finite structures, and then proves the stronger form in
which Börger, Grädel and Gurevich state it: the three classes are pairwise
recursively inseparable, so that no algorithm can even separate one class from
another.  The same section records the positive side that motivates the
question, namely that the finite model property yields a decision procedure,
and the notion of a conservative reduction, which carries the inseparability
to other formula classes.

\emph{Inputs.}  A finite relational signature \(\Lang\), that is, a finite
signature with relation symbols and no function or constant symbols;
the fixed effective coding of first-order formulas used in
\cref{sec:decision-hierarchy-01}, under which the set of codes of all
formulas and, for each finite signature, the set of codes of its sentences
are decidable; the deterministic single-tape Turing machines of
\cref{sec:turing-machines-01}, with the program codes of
\cref{sec:universal-reductions-01}; and the self-application domain set
\(K\) of \cref{sec:universal-reductions-03}.  Structures have nonempty
domains, and equality is interpreted as identity
(\cref{sec:fol-semantics-01}).  ``Recursive'' in the source's terminology
means ``decidable'' here.

\emph{Conclusions.}  For every finite relational \(\Lang\) containing a binary
relation symbol: the set of finitely satisfiable \(\Lang\)-sentences is
computably enumerable but undecidable, and the set of finitely valid
\(\Lang\)-sentences is co-computably enumerable but not computably enumerable
(\cref{thm:p02-w02-p02-036-trakhtenbrot}); the class of sentences having only
infinite models is nonempty and neither decidable nor computably enumerable,
and so is the set of finitely valid sentences that are not valid
(\cref{thm:p02-w02-p02-036-classes}); and the finitely satisfiable, the
unsatisfiable and the only infinitely satisfiable sentences are pairwise
recursively inseparable, both for \(\Lang\)-sentences and for all first-order
formulas (\cref{thm:p02-w02-p02-036-inseparable}).  A decidable class with the
finite model property has a decidable satisfiability problem
(\cref{thm:p02-w02-p02-036-fmp}), and every conservative reduction class
inherits the inseparability (\cref{thm:p02-w02-p02-036-conservative}).  All
proofs are in the annex.

\subsection{The three countermodel classes}
\label{sec:p02-w02-p02-036-classes}

\begin{hbdefinition}{def:p02-w02-p02-036-classes}{Satisfiability and countermodel classes}
Let \(X\) be a class of formulas.  Following
\cite[Sect.~2.1.2, Def.~2.1.27]{BorgerGradelGurevich1997}, write
\begin{itemize}
\item \(\mathrm{Sat}(X)\) for the set of \(\psi\in X\) that are satisfiable;
\item \(\mathrm{Val}(X)\) for the set of logically valid \(\psi\in X\);
\item \(\mathrm{Fin\text{-}sat}(X)\) for the set of \(\psi\in X\) that have a
      finite model;
\item \(\mathrm{Inf\text{-}axioms}(X)=\mathrm{Sat}(X)\setminus
      \mathrm{Fin\text{-}sat}(X)\), the \emph{infinity axioms} of \(X\): the
      satisfiable formulas all of whose models are infinite;
\item \(\mathrm{Non\text{-}sat}(X)\) for the set of unsatisfiable
      \(\psi\in X\).
\end{itemize}
A formula is satisfiable when some structure and some assignment satisfy it,
and it has a finite model when such a structure can be chosen finite.  When
\(X\) is the set \(\mathrm{FO}\) of all first-order formulas the argument is
dropped: \(\mathrm{Sat}\), \(\mathrm{Fin\text{-}sat}\), and so on.  For a
finite relational signature \(\Lang\), the main case below is
\(X=\operatorname{Sent}(\Lang)\).  A sentence \(\varphi\) is \emph{finitely
valid} when it is true in every finite structure for its signature.  A
\emph{countermodel} of \(\varphi\) is a model of \(\neg\varphi\).  As sets of
formulas these are sets of codes under the fixed effective coding of
\cref{sec:decision-hierarchy-01}; ``decidable'' and ``computably enumerable''
refer to those codes.
\index{finite satisfiability}\index{finite validity}\index{infinity axiom}
\end{hbdefinition}

For a sentence \(\varphi\) of a signature \(\Lang\), whether it is
satisfiable, or has a finite model, does not depend on the signature in which
it is read: a model for a larger signature has a reduct to \(\Lang\) with the
same domain that satisfies \(\varphi\), and an \(\Lang\)-model expands, with
the same domain, to any larger signature (\cref{lem:reduct-invariance}).
Hence \(\mathrm{Fin\text{-}sat}(\operatorname{Sent}(\Lang))=
\mathrm{Fin\text{-}sat}\cap\operatorname{Sent}(\Lang)\), and likewise for the
other classes.

The three sets \(\mathrm{Fin\text{-}sat}(X)\), \(\mathrm{Inf\text{-}axioms}(X)\)
and \(\mathrm{Non\text{-}sat}(X)\) partition \(X\): a formula either has a
finite model, or has models but only infinite ones, or has no model.  Read
through negation, the same partition classifies sentences by their
countermodels.  For \(X=\operatorname{Sent}(\Lang)\), since \(\neg\varphi\)
and \(\varphi\) have complementary classes of models,
\begin{itemize}
\item \(\varphi\) is valid exactly when \(\neg\varphi\in
      \mathrm{Non\text{-}sat}(X)\);
\item \(\varphi\) is finitely valid exactly when \(\neg\varphi\notin
      \mathrm{Fin\text{-}sat}(X)\);
\item \(\varphi\) is finitely valid but not valid exactly when
      \(\neg\varphi\in\mathrm{Inf\text{-}axioms}(X)\), that is, when
      \(\varphi\) has countermodels and every countermodel of \(\varphi\) is
      infinite.
\end{itemize}
The gap between finite and unrestricted validity is therefore exactly the
class of infinity axioms, read through negation.

\begin{hbexample}{def:p02-w02-p02-036-infinity-axiom}{The infinity axiom of the source}
Let \(K\) be a binary relation symbol.  The source's example of an infinity
axiom \cite[Sect.~2.1.2, p.~33]{BorgerGradelGurevich1997} is
\[
 \prec_K\;\defeq\;\forall x\,\exists u\,\forall y\,
 \bigl(\neg Kxx\land Kxu\land(Kyx\to Kyu)\bigr).
\]
Its intended model is the natural numbers with \(Kxu\) read as \(x<u\).  The
prefix is \(\forall\exists\forall\), and the quantifier-free part is a
conjunction of the clauses \(\neg Kxx\), \(Kxu\) and \(\neg Kyx\lor Kyu\).
Each clause has at most two literals and at most one unnegated literal, so
\(\prec_K\) lies in the class \([\forall\exists\forall]\cap\mathrm{KROM}\cap
\mathrm{HORN}\) of prenex Krom and Horn sentences with that prefix
\cite[Sect.~2.1.1, Def.~2.1.14]{BorgerGradelGurevich1997}.
Lemma~\ref{thm:p02-w02-p02-036-infinity-axiom}, which is the source's
Exercise~2.1.26, shows that
\(\prec_K\in\mathrm{Inf\text{-}axioms}(\operatorname{Sent}(\{K\}))\).
Consequently \(\neg\prec_K\) is finitely valid, is not valid, and each of its
countermodels, for instance \((\mathbb N,<)\), is infinite.  Together with
\(\exists x\,Kxx\in\mathrm{Fin\text{-}sat}\) and
\(\exists x\,(Kxx\land\neg Kxx)\in\mathrm{Non\text{-}sat}\), this shows that all
three classes are nonempty already for the signature \(\{K\}\).  The relation
symbol \(K\) of this example is unrelated to the self-application set \(K\)
of \cref{sec:universal-reductions-03}; each occurrence below is one or the
other by context.
\end{hbexample}

\begin{hblemma}{thm:p02-w02-p02-036-infinity-axiom}{proved=annex:p02-w02-p02-036-infinity-axiom}{The sentence \(\prec_K\) has only infinite models}
The structure \((\mathbb N,<)\), with \(<\) interpreting \(K\), is a model of
\(\prec_K\), and \(\prec_K\) has no finite model.
\end{hblemma}

\subsection{Finite structures as finite data}
\label{sec:p02-w02-p02-036-finite-data}

\begin{hbconvention}{def:p02-w02-p02-036-finite-data}{Finite structures as machine inputs}
Fix a finite relational signature \(\Lang\) with its relation symbols listed
as \(R_1,\ldots,R_r\), where \(R_i\) has arity \(m_i\).  Every finite
\(\Lang\)-structure is isomorphic to one whose domain is
\(\{0,\ldots,n-1\}\) for some \(n\geq1\), and isomorphic structures satisfy
the same sentences (\cref{lem:iso-invariance}).  A structure with domain
\(\{0,\ldots,n-1\}\) is given to a machine as the number \(n\) together with,
for each \(i\), the finite list of the \(m_i\)-tuples in \(R_i\), each relation
being a subset of the finite set \(\{0,\ldots,n-1\}^{m_i}\).  This datum is a
finite word over a fixed finite alphabet, and for each \(n\) there are only
finitely many such data.  One concrete format, the one Fagin uses, writes \(n\)
in binary followed by the characteristic words of the relations in
lexicographic order; it is fixed in
Convention~\ref{def:p02-w02-p02-036-fagin-encoding}.  Any format from which
\(n\) and the relations are computably recoverable serves for the results of
this section.  When the signature also has function and constant symbols, a
function symbol of arity \(m\) is given by its table, the list of its \(n^m\)
values, and a constant symbol by its value.
\end{hbconvention}

\begin{hblemma}{thm:p02-w02-p02-036-finsat-ce}{proved=annex:p02-w02-p02-036-finsat-ce}{Finite model checking and enumerability}
Let \(\Lang\) be a finite relational signature and \(X=\operatorname{Sent}(\Lang)\).
\begin{enumerate}
\item There is an algorithm that, given a finite \(\Lang\)-structure with domain
      \(\{0,\ldots,n-1\}\) in the format of
      Convention~\ref{def:p02-w02-p02-036-finite-data} and a sentence
      \(\varphi\in X\), decides whether the structure satisfies \(\varphi\).
\item \(\mathrm{Fin\text{-}sat}(X)\) is computably enumerable.
\item \(\mathrm{Non\text{-}sat}(X)\) is computably enumerable.
\item The set of finitely valid sentences in \(X\) is co-computably enumerable:
      its complement in \(X\) is computably enumerable.
\end{enumerate}
\end{hblemma}

The enumeration in item~2 searches through finite structures; it halts on a
sentence exactly when a finite model turns up.  It gives no information when
the sentence has no finite model, which is the case Trakhtenbrot's theorem
shows cannot be detected effectively.  The same two facts, enumerability of
the finite structures up to isomorphism and decidability of finite model
checking, are the source's ground for the following decision method
\cite[Sect.~2.1.2, p.~33]{BorgerGradelGurevich1997}.  A class \(X\) of
formulas has the \emph{finite model property} when every satisfiable member of
\(X\) has a finite model, that is, when \(\mathrm{Sat}(X)=
\mathrm{Fin\text{-}sat}(X)\).

\begin{hbtheorem}{thm:p02-w02-p02-036-fmp}{proved=annex:p02-w02-p02-036-fmp}{The finite model property gives decidability}
Let \(\Lang\) be a finite signature, possibly with function and constant
symbols, and let \(X\) be a decidable set of \(\Lang\)-formulas with the
finite model property.  Then \(\mathrm{Sat}(X)\) is decidable.
\end{hbtheorem}

The search through finite structures supplies one half of the decision
procedure: it confirms the satisfiable members.  The other half, confirming
the unsatisfiable members, comes from the completeness theorem, through the
enumeration of the valid sentences.  The finite model property is what makes
the two searches together cover every member of \(X\).  Trakhtenbrot's theorem
shows that the first search alone cannot be completed to a decision procedure
for all sentences of a signature with a binary relation symbol.

\subsection{Accepting computations as finite structures}
\label{sec:p02-w02-p02-036-run-sentence}

Let \(\mathcal T=(Q,\mathcal A_{\mathrm{in}},\mathcal A_{\mathrm{tape}},
\delta,q_0,q_{\mathrm{acc}},q_{\mathrm{rej}})\) be a Turing machine in the
sense of \cref{sec:turing-machines-01}, with its two-way infinite tape.  The
structure built below is a finite table whose rows are moments of time and whose
columns are tape cells; one linear order serves for both.

\begin{hbdefinition}{def:p02-w02-p02-036-run-sentence}{The run sentence of a machine and an input}
The \emph{run signature} \(\Lang_{\mathcal T}\) consists of a binary relation
symbol \(<\), a binary symbol \(\mathsf{Hd}\), a binary symbol
\(\mathsf{Cl}_a\) for each \(a\in\mathcal A_{\mathrm{tape}}\), and a unary
symbol \(\mathsf{St}_q\) for each \(q\in Q\).  The intended reading of
\(\mathsf{St}_q(t)\) is ``at time \(t\) the state is \(q\)'', of
\(\mathsf{Hd}(t,p)\) ``at time \(t\) the head scans cell \(p\)'', and of
\(\mathsf{Cl}_a(t,p)\) ``at time \(t\) cell \(p\) holds \(a\)''.  Use the
abbreviations
\[
 \mathsf{Nx}(x,y)\defeq x<y\land\neg\exists z\,(x<z\land z<y),
 \qquad
 \mathsf{Fst}(x)\defeq\neg\exists z\,z<x ,
\]
and \(\mathsf{Mv}_R(p,p')\defeq\mathsf{Nx}(p,p')\),
\(\mathsf{Mv}_L(p,p')\defeq\mathsf{Nx}(p',p)\).  For an input word
\(w=a_0a_1\cdots a_{k-1}\in\mathcal A_{\mathrm{in}}^*\) let
\(\psi_{\mathcal T,w}\) be the conjunction of the following sentences.
\begin{enumerate}
\item \emph{Order.} \(\forall x\,\neg x<x\),
  \(\forall x\,\forall y\,\forall z\,((x<y\land y<z)\to x<z)\) and
  \(\forall x\,\forall y\,(x<y\lor x=y\lor y<x)\).
\item \emph{Determinacy.} For every \(t\): exactly one \(q\in Q\) has
  \(\mathsf{St}_q(t)\); exactly one \(p\) has \(\mathsf{Hd}(t,p)\); and for
  every \(p\), exactly one \(a\in\mathcal A_{\mathrm{tape}}\) has
  \(\mathsf{Cl}_a(t,p)\).  These are written with finite disjunctions over
  \(Q\) and \(\mathcal A_{\mathrm{tape}}\), pairwise exclusions, and
  \(\forall t\,\exists p\,\mathsf{Hd}(t,p)\land\forall t\,\forall p\,\forall p'\,
  ((\mathsf{Hd}(t,p)\land\mathsf{Hd}(t,p'))\to p=p')\).
\item \emph{Start.} \(\forall t\,(\mathsf{Fst}(t)\to\mathsf{St}_{q_0}(t))\) and
  \(\forall t\,\forall p\,((\mathsf{Fst}(t)\land\mathsf{Hd}(t,p))\to
  \mathsf{In}_w(t,p))\), where \(\mathsf{In}_w(t,p)\) is
  \(\forall u\,\mathsf{Cl}_{\sqcup}(t,u)\) when \(k=0\) and otherwise
  \[
  \begin{aligned}
   \exists y_0\cdots\exists y_{k-1}\Bigl(&y_0=p\land
   \bigwedge_{i<k-1}\mathsf{Nx}(y_i,y_{i+1})\land
   \bigwedge_{i<k}\mathsf{Cl}_{a_i}(t,y_i)\\
   &\land\forall u\,\bigl(\textstyle\bigwedge_{i<k}u\neq y_i\to
   \mathsf{Cl}_{\sqcup}(t,u)\bigr)\Bigr).
  \end{aligned}
  \]
\item \emph{Steps.} For each \(q\notin\{q_{\mathrm{acc}},q_{\mathrm{rej}}\}\)
  and \(a\in\mathcal A_{\mathrm{tape}}\) with \(\delta(q,a)=(q',b,d)\), the
  sentence
  \[
  \begin{aligned}
   \forall t\,\forall p\,\Bigl(&(\mathsf{St}_q(t)\land\mathsf{Hd}(t,p)\land
     \mathsf{Cl}_a(t,p))\to
     \exists t'\,\exists p'\,\bigl(\mathsf{Nx}(t,t')\land
     \mathsf{Mv}_d(p,p')\\
   &\land\mathsf{St}_{q'}(t')\land\mathsf{Hd}(t',p')\land\mathsf{Cl}_b(t',p)\\
   &\land\forall u\,\bigl(u\neq p\to\textstyle\bigwedge_{c\in\mathcal A_{\mathrm{tape}}}
     (\mathsf{Cl}_c(t,u)\to\mathsf{Cl}_c(t',u))\bigr)\bigr)\Bigr).
  \end{aligned}
  \]
\item \emph{Acceptance.} \(\exists t\,\mathsf{St}_{q_{\mathrm{acc}}}(t)\land
  \neg\exists t\,\mathsf{St}_{q_{\mathrm{rej}}}(t)\).
\end{enumerate}
\index{run sentence}
\end{hbdefinition}

The step clause forces a next moment and a neighbouring cell to exist
whenever the machine is not halted, so a finite model can only describe a
computation that halts inside it.  The acceptance clause forbids the rejecting
state anywhere, so that the first halting moment must be accepting.  No
clause says where the head starts: the tape is described relative to the
initial head cell, and the step relation is the same at every cell.  The run
sentence controls finite models only; it can have infinite models, for
example orders without a least element, that describe no computation, so it
does not by itself separate the unsatisfiable sentences.  That is the job of
the refined formula of \cref{sec:p02-w02-p02-036-refined}.

\begin{hblemma}{thm:p02-w02-p02-036-run-sentence}{proved=annex:p02-w02-p02-036-run-sentence}{Accepting runs and finite models}
For every Turing machine \(\mathcal T\) and every
\(w\in\mathcal A_{\mathrm{in}}^*\), the computation of \(\mathcal T\) on \(w\)
reaches \(q_{\mathrm{acc}}\) if and only if \(\psi_{\mathcal T,w}\) has a finite
model.  If that computation reaches \(q_{\mathrm{acc}}\) after \(s\) steps, then
\(\psi_{\mathcal T,w}\) has a model with exactly \(2s+\lvert w\rvert+1\)
elements.  For fixed \(\mathcal T\), the map \(w\mapsto\psi_{\mathcal T,w}\) is
computable.
\end{hblemma}

\subsection{Reduction to one binary relation}
\label{sec:p02-w02-p02-036-one-binary}

The run signature has many symbols.  The next construction replaces them by a
single binary relation \(E\): domain elements are the vertices with a loop,
and each fact \(R_j(a,b)\) becomes a path of loop-free vertices from \(a\) to
\(b\) whose length records the index \(j\).

\begin{hbdefinition}{def:p02-w02-p02-036-gadget}{The path translation}
Let \(\Lang\) consist of binary relation symbols \(R_1,\ldots,R_l\) and unary
relation symbols \(P_1,\ldots,P_m\), and let \(E\) be a binary relation
symbol.  Give \(R_j\) the index \(j\) and \(P_i\) the index \(l+i\).  For
\(j\geq1\) and variables \(x,y\), let \(z_1,\ldots,z_j\) be the first \(j\)
variables, in a fixed enumeration of the variables, that differ from \(x\) and
\(y\), and put
\[
 \pi_j(x,y)\defeq\exists z_1\cdots\exists z_j\Bigl(Exz_1\land
 \bigwedge_{1\leq i<j}Ez_iz_{i+1}\land Ez_jy\land
 \bigwedge_{1\leq i\leq j}\neg Ez_iz_i\Bigr).
\]
The translation \(\varphi\mapsto\varphi^{I}\) from \(\Lang\)-formulas to
\(\{E\}\)-formulas sends \(R_j(x,y)\) to \(\pi_j(x,y)\) and \(P_i(x)\) to
\(\pi_{l+i}(x,x)\), leaves \(x=y\) and \(\bottom\) unchanged, commutes with
\(\neg,\land,\lor,\to\), and relativizes quantifiers to looped vertices:
\[
 (\forall x\,\varphi)^{I}\defeq\forall x\,(Exx\to\varphi^{I}),\qquad
 (\exists x\,\varphi)^{I}\defeq\exists x\,(Exx\land\varphi^{I}).
\]
For a sentence \(\varphi\) put
\(\varphi^{\sharp}\defeq\exists x\,Exx\land\varphi^{I}\).

For an \(\{E\}\)-structure \(\mathcal G\) in which some vertex has a loop, the
\emph{decoded structure} \(D(\mathcal G)\) is the \(\Lang\)-structure whose
domain is the set of looped vertices, with
\(R_j^{D(\mathcal G)}=\{(a,b):\mathcal G\entails\pi_j[a,b]\}\) and
\(P_i^{D(\mathcal G)}=\{a:\mathcal G\entails\pi_{l+i}[a,a]\}\), \(a,b\) ranging
over looped vertices.  For an \(\Lang\)-structure \(\mathcal M\), the
\emph{path graph} \(G(\mathcal M)\) has as vertices the elements of
\(\operatorname{dom}\mathcal M\) and, for each \emph{fact}
\(f=(a,b,j)\), meaning \((a,b)\in R_j^{\mathcal M}\) with \(j\leq l\) or
\(a=b\in P_i^{\mathcal M}\) with \(j=l+i\), fresh vertices
\(c^f_1,\ldots,c^f_j\).  Its edges are the loops \((a,a)\) for
\(a\in\operatorname{dom}\mathcal M\) and, for each fact \(f=(a,b,j)\), the
edges \(a\to c^f_1\), \(c^f_i\to c^f_{i+1}\) for \(1\leq i<j\), and
\(c^f_j\to b\).
\end{hbdefinition}

\begin{hblemma}{thm:p02-w02-p02-036-gadget}{proved=annex:p02-w02-p02-036-gadget}{One binary relation preserves satisfiability and finite satisfiability}
In the setting of Definition~\ref{def:p02-w02-p02-036-gadget}:
\begin{enumerate}
\item For every \(\{E\}\)-structure \(\mathcal G\) with a looped vertex, every
      \(\Lang\)-formula \(\varphi\) and every assignment \(s\) with values
      among the looped vertices,
      \(\mathcal G,s\entails\varphi^{I}\) if and only if
      \(D(\mathcal G),s\entails\varphi\).
\item For every \(\Lang\)-structure \(\mathcal M\), \(D(G(\mathcal M))=\mathcal M\);
      and \(G(\mathcal M)\) is finite when \(\mathcal M\) is finite.
\item For every \(\varphi\in\operatorname{Sent}(\Lang)\),
      \(\varphi\in\mathrm{Fin\text{-}sat}(\operatorname{Sent}(\Lang))\) if and
      only if \(\varphi^{\sharp}\in\mathrm{Fin\text{-}sat}
      (\operatorname{Sent}(\{E\}))\), and \(\varphi\mapsto\varphi^{\sharp}\) is
      computable.
\item For every \(\varphi\in\operatorname{Sent}(\Lang)\), the sentence
      \(\varphi\) is satisfiable if and only if \(\varphi^{\sharp}\) is
      satisfiable.  Consequently \(\varphi\) and \(\varphi^{\sharp}\) lie in corresponding
      classes among \(\mathrm{Fin\text{-}sat}\), \(\mathrm{Inf\text{-}axioms}\)
      and \(\mathrm{Non\text{-}sat}\).
\end{enumerate}
\end{hblemma}

\begin{hblemma}{thm:p02-w02-p02-036-transfer}{proved=annex:p02-w02-p02-036-transfer}{Transfer to any signature with a binary symbol}
Let \(\Lang'\) be a finite relational signature containing a binary relation
symbol \(R\).  For \(\chi\in\operatorname{Sent}(\{E\})\) let \(\chi[R/E]\) be
the result of replacing every occurrence of \(E\) by \(R\).  Then
\(\chi\in\mathrm{Fin\text{-}sat}(\operatorname{Sent}(\{E\}))\) if and only if
\(\chi[R/E]\in\mathrm{Fin\text{-}sat}(\operatorname{Sent}(\Lang'))\), and
\(\chi\in\mathrm{Sat}(\operatorname{Sent}(\{E\}))\) if and only if
\(\chi[R/E]\in\mathrm{Sat}(\operatorname{Sent}(\Lang'))\).
\end{hblemma}

In the terminology of Definition~\ref{def:p02-w02-p02-036-conservative}
below, \(\varphi\mapsto\varphi^{\sharp}\) and \(\chi\mapsto\chi[R/E]\) are
conservative reductions: they preserve and reflect both satisfiability and
finite satisfiability.

\subsection{Trakhtenbrot's theorem and the classes}
\label{sec:p02-w02-p02-036-conclusions}

\begin{hbcorollary}{thm:p02-w02-p02-036-trakhtenbrot}{proved=annex:p02-w02-p02-036-trakhtenbrot}{Proof of Trakhtenbrot's theorem}
Let \(\Lang\) be a finite relational signature containing a binary relation
symbol and \(X=\operatorname{Sent}(\Lang)\).  Then
\(\mathrm{Fin\text{-}sat}(X)\) is computably enumerable and not decidable, and
the set of finitely valid sentences in \(X\) is co-computably enumerable and
not computably enumerable.  This is the statement of
\cref{thm:trakhtenbrot}.
\end{hbcorollary}

\begin{hbtheorem}{thm:p02-w02-p02-036-classes}{proved=annex:p02-w02-p02-036-classes}{Effective standing of the countermodel classes}
Let \(\Lang\) be a finite relational signature containing a binary relation
symbol and \(X=\operatorname{Sent}(\Lang)\).
\begin{enumerate}
\item \(\mathrm{Fin\text{-}sat}(X)\), \(\mathrm{Non\text{-}sat}(X)\) and their
      union \(X\setminus\mathrm{Inf\text{-}axioms}(X)\) are computably
      enumerable.
\item \(\mathrm{Inf\text{-}axioms}(X)\) is nonempty, not decidable, and not
      computably enumerable.
\item The set of sentences in \(X\) that are finitely valid but not valid is
      nonempty and not computably enumerable.  Each of its members has a
      countermodel, and all of its countermodels are infinite.
\end{enumerate}
\end{hbtheorem}

\subsection{Recursive inseparability and three halting problems}
\label{sec:p02-w02-p02-036-inseparability}

Undecidability of a set says that no algorithm draws its boundary exactly.
The source's form of Trakhtenbrot's theorem says more: no algorithm draws
any boundary that keeps two of the three classes apart.

\begin{hbdefinition}{def:p02-w02-p02-036-inseparable}{Recursive inseparability}
Two disjoint sets \(X,Y\subseteq\mathbb N\) are called \emph{recursively
inseparable} if there is no recursive set \(R\) such that \(X\subseteq R\) and
\(R\cap Y=\varnothing\) \cite[Sect.~2.1.2, Def.~2.1.28]{BorgerGradelGurevich1997}.
A set \(R\) with these two properties \emph{separates} \(X\) from \(Y\).
Sets of formulas are read as sets of codes.
\index{recursive inseparability}
\end{hbdefinition}

If \(R\) separates \(X\) from \(Y\), then the decidable set
\(\mathbb N\setminus R\) separates \(Y\) from \(X\); so the relation is
symmetric.  If \(X\cup Y=\mathbb N\) and \(X\) is not decidable, then \(X\) and
\(Y\) are recursively inseparable, since the only set separating \(X\) from
\(Y\) is \(X\) itself.  Recursive inseparability of a pair therefore
strengthens the undecidability of each member.

The machines in the next definition are those of
\cref{sec:turing-machines-01}, started on the empty input, and they are named
by the program codes of \cref{sec:universal-reductions-01}: a number \(e\) is a
\emph{valid code} when it codes a normalized machine description, and then
\(M_e\) denotes the decoded machine.  The set \(V\) of valid codes is
decidable (\cref{sec:universal-reductions-01}).

\begin{hbdefinition}{def:p02-w02-p02-036-halting}{The halting problems \(H_1\), \(H_2\) and \(H_\infty\)}
For a valid code \(e\), consider the computation of \(M_e\) on the empty
input.  Put
\[
\begin{aligned}
 H_1&\defeq\{e\in V:\text{the computation of }M_e\text{ on the empty input
      reaches }q_{\mathrm{rej}}\},\\
 H_2&\defeq\{e\in V:\text{the computation of }M_e\text{ on the empty input
      reaches }q_{\mathrm{acc}}\},\\
 H_\infty&\defeq\{e\in V:\text{the computation of }M_e\text{ on the empty
      input never halts}\}.
\end{aligned}
\]
These three sets partition \(V\).  The source defines \(H_1\), \(H_2\) and the
non-halting set for its Turing machine programs with exactly two stop states
\(1\) and \(2\), started in state \(0\) at the left end of the empty one-way
infinite tape \cite[Sect.~2.1.2, p.~34]{BorgerGradelGurevich1997}.  Here the
rejecting state plays the role of the source's stop state \(1\) and the
accepting state the role of its stop state \(2\); the tape is the two-way tape
of \cref{sec:turing-machines-01}.
\index{halting problem!recursively inseparable}
\end{hbdefinition}

\begin{hblemma}{thm:p02-w02-p02-036-halting}{proved=annex:p02-w02-p02-036-halting}{The three halting problems are pairwise recursively inseparable}
The sets \(H_1\), \(H_2\) and \(H_\infty\) are pairwise recursively
inseparable.
\end{hblemma}

The source takes this lemma from the recursion-theory literature and also
records that the three sets are recursively enumerable
\cite[Sect.~2.1.2, p.~34]{BorgerGradelGurevich1997}; enumerability plays no
part in the arguments below, and the annex proves the inseparability from the
definition of a decider.

\subsection{The refined reduction formula}
\label{sec:p02-w02-p02-036-refined}

The source obtains Trakhtenbrot's theorem from the reduction formula \(\rho_M\)
of its proof of the Church--Turing theorem.  That formula is the universal
closure of a conjunction \(\mathrm{STEP}_M\land\mathrm{START}\land
\mathrm{NONSTOP}\land\mathrm{ORDER}\), in which \(\mathrm{ORDER}\) makes \(K\)
a strict linear order with successor relation \(S\) and also requires every
element to have a successor, so that every model is infinite
\cite[Sect.~2.1.1, pp.~20--22]{BorgerGradelGurevich1997}.  For the refined
property the source replaces \(\mathrm{ORDER}\) by a relativization
\(\prec_{K,M}\) of the infinity axiom \(\prec_K\) to the cells visited by the
computation, so that a halting computation has a finite model while a
non-halting one has only infinite models
\cite[Sect.~2.1.2, pp.~34--35]{BorgerGradelGurevich1997}.  The definition below
carries this out in a relational signature for the machines of
\cref{sec:turing-machines-01}.

\begin{hbdefinition}{def:p02-w02-p02-036-refined}{The refined reduction formula \(\rho_{M,\prec}\)}
Let \(M=(Q,\mathcal A_{\mathrm{in}},\mathcal A_{\mathrm{tape}},\delta,q_0,
q_{\mathrm{acc}},q_{\mathrm{rej}})\) be a Turing machine.  Its signature
\(\Lang_M^{\prec}\) consists of the binary relation symbols \(K\), \(S\),
\(H\) and \(T_a\) for \(a\in\mathcal A_{\mathrm{tape}}\), and the unary
relation symbols \(Z\) and \(I_q\) for \(q\in Q\).  The intended readings,
over the integers, are those of the source
\cite[Sect.~2.1.1, p.~20]{BorgerGradelGurevich1997}: \(H(t,x)\), the head
position at time \(t\) is cell \(x\); \(T_a(t,x)\), the tape at time \(t\)
contains \(a\) in cell \(x\); \(I_q(t)\), at time \(t\) the machine is in state
\(q\); \(S(t,t')\), \(t'\) is the successor of \(t\); \(K(x,y)\), \(x<y\); and
\(Z(x)\), \(x=0\).  Write \(\mathrm{Act}(x)\) for the disjunction of the
\(I_q(x)\) with \(q\notin\{q_{\mathrm{acc}},q_{\mathrm{rej}}\}\), that is,
``at time \(x\) the machine is not halted''.  The sentence
\(\rho_{M,\prec}\) is the conjunction of the following sentences.
\begin{enumerate}
\item \(\mathrm{ORDER}^{-}\): the source's \(\mathrm{ORDER}\) without its
  successor-existence clause,
  \[
  \begin{gathered}
   \forall x\,\neg K(x,x),\qquad
   \forall x\,\forall y\,\bigl(K(x,y)\lor x=y\lor K(y,x)\bigr),\\
   \forall x\,\forall y\,\forall z\,\bigl((K(x,y)\land K(y,z))\to K(x,z)\bigr),\\
   \forall x\,\forall y\,\bigl(S(x,y)\leftrightarrow(K(x,y)\land
   \neg\exists z\,(K(x,z)\land K(z,y)))\bigr),
  \end{gathered}
  \]
  where \(\leftrightarrow\) abbreviates the conjunction of the two
  implications.
\item \(\mathrm{START}\): \(\exists x\,Z(x)\land\forall x\,\bigl(Z(x)\to
  (I_{q_0}(x)\land H(x,x)\land\forall y\,T_{\sqcup}(x,y))\bigr)\).
\item \(\mathrm{STEP}_M\), the universal closures of the following
  implications.  For each \(q\notin\{q_{\mathrm{acc}},q_{\mathrm{rej}}\}\) and
  \(a\in\mathcal A_{\mathrm{tape}}\) with \(\delta(q,a)=(q',b,R)\):
  \[
  \begin{aligned}
   &\bigl(I_q(t)\land H(t,x)\land T_a(t,x)\land S(t,t')\land S(x,x')\bigr)\\
   &\qquad\to\bigl(I_{q'}(t')\land H(t',x')\land T_b(t',x)\bigr);
  \end{aligned}
  \]
  with \(\delta(q,a)=(q',b,L)\):
  \[
  \begin{aligned}
   &\bigl(I_q(t)\land H(t,x')\land T_a(t,x')\land S(t,t')\land S(x,x')\bigr)\\
   &\qquad\to\bigl(I_{q'}(t')\land H(t',x)\land T_b(t',x')\bigr);
  \end{aligned}
  \]
  and, for each \(a\in\mathcal A_{\mathrm{tape}}\), no change outside the
  head cell:
  \[
   \bigl(H(t,x)\land y\neq x\land T_a(t,y)\land S(t,t')\bigr)\to T_a(t',y).
  \]
\item \(\mathrm{NONSTOP}\): \(\forall t\,\neg I_{q_{\mathrm{rej}}}(t)\).
\item \(\prec_{K,M}\), the relativized infinity axiom:
  \[
   \forall x\,\exists u\,\exists v\,\forall y\,
   \Bigl(\bigl((H(y,x)\lor\mathrm{Act}(x))\to S(x,u)\bigr)\land
   \bigl(H(y,x)\to S(v,x)\bigr)\Bigr).
  \]
\end{enumerate}
Finally put \(\rho^{\infty}_{M}\defeq\rho_{M,\prec}\land
\forall x\,\exists y\,S(x,y)\), which restores the unrestricted successor
clause of \(\mathrm{ORDER}\).
\index{reduction formula!refined}
\end{hbdefinition}

The step clauses are the source's formulas (2.1)--(2.3) for a deterministic
machine \cite[Sect.~2.1.1, pp.~20--21]{BorgerGradelGurevich1997}, and
\(\mathrm{START}\) and \(\mathrm{NONSTOP}\) are the source's, with the constant
\(0\) replaced by the unary predicate \(Z\) and the halting state \(1\) by
\(q_{\mathrm{rej}}\).  Like the source's, these formulas are frugal: they do
not say that a moment has only one state or a cell only one letter.  The
relativized axiom \(\prec_{K,M}\) asserts successors only where the
computation needs them: a successor moment after every moment at which the
machine is not halted, and both neighbours of every visited cell.  Together
with \(\mathrm{ORDER}^{-}\), under which \(S\) is contained in \(K\), it
yields for each such element a \(K\)-larger one, which is what the source's
display \(\forall x\,\exists u\,\forall y\,(\neg Kxx\land(Hyx\to Kxu)\land
(Kyx\to Kyu))\) asserts for visited cells
\cite[Sect.~2.1.2, p.~35]{BorgerGradelGurevich1997}.  The source reads the
successor as a function \(x'\), so that successors always exist; in the
relational version they must be asserted, and the relativization asserts
them for active moments as well as visited cells.  For the same reason the
source's sharpening of the stop condition, by which a non-halting computation
never repeats a configuration, is not needed here: the active moments of a
non-halting computation already form an infinite \(K\)-chain.

\begin{hblemma}{thm:p02-w02-p02-036-simulation}{proved=annex:p02-w02-p02-036-simulation}{Simulation lemma}
Let \(M\) be a Turing machine, let \(c_0,c_1,\ldots\) be its computation on
the empty input, with \(c_t=(q_t,B_t,h_t)\), and let \(\mathfrak A\) be a
model of \(\rho_{M,\prec}\) with an element \(z\in Z^{\mathfrak A}\).  Then there
are elements \(\tau_t\), one for each \(t\) such that \(c_t\) exists, and an
injective map \(\sigma\) from an interval \(I\subseteq\mathbb Z\) containing
\(0\) into the domain of \(\mathfrak A\), such that \(\tau_0=\sigma_0=z\),
\(K^{\mathfrak A}(\tau_s,\tau_t)\) whenever \(s<t\), and for every \(t\) such
that \(c_t\) exists:
\begin{enumerate}
\item \(\mathfrak A\entails I_{q_t}[\tau_t]\);
\item \(h_t\in I\) and \(\mathfrak A\entails H[\tau_t,\sigma_{h_t}]\);
\item \(\mathfrak A\entails T_{B_t(x)}[\tau_t,\sigma_x]\) for every \(x\in I\),
      and \(\mathfrak A\entails T_{\sqcup}[\tau_t,y]\) for every element \(y\)
      outside the image of \(\sigma\).
\end{enumerate}
In particular the elements \(\tau_t\) are pairwise distinct.
\end{hblemma}

\begin{hbtheorem}{thm:p02-w02-p02-036-refined}{proved=annex:p02-w02-p02-036-refined}{Refined reduction property}
For every Turing machine \(M\), started on the empty input:
\begin{enumerate}
\item the computation reaches \(q_{\mathrm{rej}}\) if and only if
      \(\rho_{M,\prec}\) is contradictory;
\item the computation reaches \(q_{\mathrm{acc}}\) if and only if
      \(\rho_{M,\prec}\) is finitely satisfiable; if it does so after \(k\)
      steps, \(\rho_{M,\prec}\) has a model with exactly \(2k+3\) elements;
\item the computation never halts if and only if \(\rho_{M,\prec}\) is an
      infinity axiom;
\item the computation reaches \(q_{\mathrm{rej}}\) if and only if
      \(\rho^{\infty}_{M}\) is contradictory, and otherwise
      \(\rho^{\infty}_{M}\) is an infinity axiom.
\end{enumerate}
The maps \(M\mapsto\rho_{M,\prec}\) and \(M\mapsto\rho^{\infty}_{M}\) are
computable from the transition table.
\end{hbtheorem}

Items 1 and~2 are the source's reduction property for \(\rho_{M,\prec}\), with
\(H_1\) and \(H_2\) read as in
Definition~\ref{def:p02-w02-p02-036-halting}
\cite[Sect.~2.1.2, p.~34]{BorgerGradelGurevich1997}.  Item~4 is the reduction
used for the source's Exercise~2.1.29.

\subsection{Pairwise recursive inseparability}
\label{sec:p02-w02-p02-036-trakhtenbrot-inseparable}

\begin{hbtheorem}{thm:p02-w02-p02-036-inseparable}{proved=annex:p02-w02-p02-036-inseparable}{Trakhtenbrot's theorem, inseparable form}
Let \(\Lang\) be a finite relational signature containing a binary relation
symbol and \(X=\operatorname{Sent}(\Lang)\).  The sets
\(\mathrm{Fin\text{-}sat}(X)\), \(\mathrm{Non\text{-}sat}(X)\) and
\(\mathrm{Inf\text{-}axioms}(X)\) are pairwise recursively inseparable.
Consequently the sets of all finitely satisfiable, all unsatisfiable and all
only infinitely satisfiable first-order formulas, \(\mathrm{Fin\text{-}sat}\),
\(\mathrm{Non\text{-}sat}\) and \(\mathrm{Inf\text{-}axioms}\), are pairwise
recursively inseparable \cite[Sect.~2.1.2, Thm.~2.1.30]{BorgerGradelGurevich1997}.
\end{hbtheorem}

\begin{hbcorollary}{thm:p02-w02-p02-036-exercise-2-1-29}{proved=annex:p02-w02-p02-036-exercise-2-1-29}{Infinity axioms and unsatisfiable formulas}
The classes \(\mathrm{Inf\text{-}axioms}\) and \(\mathrm{Non\text{-}sat}\) are
recursively inseparable.  So are the two sets
\(\mathrm{Inf\text{-}axioms}(X)\) and \(\mathrm{Non\text{-}sat}(X)\), for \(X\)
as in
Theorem~\ref{thm:p02-w02-p02-036-inseparable}
\cite[Sect.~2.1.2, Exercise~2.1.29]{BorgerGradelGurevich1997}.
\end{hbcorollary}

This pair is one of the three pairs of
Theorem~\ref{thm:p02-w02-p02-036-inseparable}; the annex proves it a second way,
by the source's route through the formula \(\rho^{\infty}_{M}\), whose models
are all infinite, and the undecidability of \(H_1\).

Recursive inseparability contains the earlier conclusions:
\begin{itemize}
\item a decidable \(\mathrm{Fin\text{-}sat}(X)\) would separate itself from
      \(\mathrm{Non\text{-}sat}(X)\);
\item a decidable \(\mathrm{Inf\text{-}axioms}(X)\) would separate itself from
      \(\mathrm{Non\text{-}sat}(X)\).
\end{itemize}
It adds that no decidable approximation of finite satisfiability can be correct on
both the finitely satisfiable and the unsatisfiable sentences, whatever it
does on the infinity axioms.

\subsection{Conservative reductions}
\label{sec:p02-w02-p02-036-conservative}

\begin{hbdefinition}{def:p02-w02-p02-036-conservative}{Conservative reduction}
Let \(X\) and \(Y\) be recursive classes of formulas.  A \emph{conservative
reduction} from \(X\) to \(Y\) is a recursive function \(g:X\to Y\) such that
for all \(\psi\in X\):
\begin{enumerate}
\item \(\psi\) is satisfiable if and only if \(g(\psi)\) is satisfiable;
\item \(\psi\) is finitely satisfiable if and only if \(g(\psi)\) is finitely
      satisfiable.
\end{enumerate}
A formula class \(X\) is a \emph{conservative reduction class} if there exists
a conservative reduction from the set of all first-order formulas to \(X\)
\cite[Sect.~2.1.2, Def.~2.1.35]{BorgerGradelGurevich1997}.  Here ``recursive
function'' means a function on codes computed by a Turing machine that halts
on every member of \(X\).
\index{conservative reduction}
\end{hbdefinition}

A conservative reduction maps each of \(\mathrm{Fin\text{-}sat}(X)\),
\(\mathrm{Inf\text{-}axioms}(X)\) and \(\mathrm{Non\text{-}sat}(X)\) into the
corresponding class of \(Y\), and it maps no member of one class into another
class, since the class of a formula is fixed by whether it is satisfiable and
whether it is finitely satisfiable.

\begin{hbtheorem}{thm:p02-w02-p02-036-conservative}{proved=annex:p02-w02-p02-036-conservative}{Conservative reduction classes inherit inseparability}
Let \(X\) be a conservative reduction class.  Then the sets
\(\mathrm{Non\text{-}sat}(X)\), \(\mathrm{Fin\text{-}sat}(X)\) and
\(\mathrm{Inf\text{-}axioms}(X)\) are pairwise recursively inseparable
\cite[Sect.~2.1.2, Exercise~2.1.36]{BorgerGradelGurevich1997}.
\end{hbtheorem}

The source motivates the definition by a sharpening of its register-machine
reduction class, and lists two prefix classes of Krom and Horn formulas as
conservative reduction classes in Exercise~2.1.37
\cite[Sect.~2.1.2, pp.~36--37]{BorgerGradelGurevich1997}; neither of those
results is proved in this section.

\subsection{Examples and limits}
\label{sec:p02-w02-p02-036-examples}

\begin{hbexample}{def:p02-w02-p02-036-worked}{A run sentence, a path graph and the refined formula}
\emph{A run sentence.}  Let \(\mathcal T\) have
\(\delta(q_0,a)=(q_{\mathrm{acc}},1,R)\) for every tape symbol \(a\), and take
the empty input.  The computation accepts after \(s=1\) step, so the proof of
Lemma~\ref{thm:p02-w02-p02-036-run-sentence} builds a model of
\(\psi_{\mathcal T,\varepsilon}\) on \(\{0,1,2\}\) with the usual order.
Real cell \(i\) is element \(i+1\), so the head starts at element \(1\).
The moments are \(0\), \(1\) and the padding moment \(2\), which repeats
moment \(1\):
\(\mathsf{St}_{q_0}=\{0\}\), \(\mathsf{St}_{q_{\mathrm{acc}}}=\{1,2\}\),
\(\mathsf{Hd}=\{(0,1),(1,2),(2,2)\}\), \(\mathsf{Cl}_1=\{(1,1),(2,1)\}\), and
\(\mathsf{Cl}_{\sqcup}\) holds at every other pair.  The only step clause with
a true antecedent is the one for \((q_0,\sqcup)\) at \(t=0\), \(p=1\); it is
witnessed by \(t'=1\), \(p'=2\).  If instead
\(\delta(q_0,a)=(q_0,a,R)\) for every \(a\), the computation never halts, and by
the lemma \(\psi_{\mathcal T,\varepsilon}\) has no finite model: in any finite
model, the step clause would need a cell to the right of the last element.

\emph{A path graph.}  Let \(\Lang=\{R_1,P_1\}\) with \(R_1\) binary and
\(P_1\) unary, so \(l=m=1\) and \(P_1\) has index \(2\).  Let \(\mathcal M\)
have domain \(\{a,b\}\), \(R_1^{\mathcal M}=\{(a,b)\}\) and
\(P_1^{\mathcal M}=\{b\}\).  Then \(G(\mathcal M)\) has vertices
\(a,b,c,d_1,d_2\) and edges \(a\to a\), \(b\to b\), \(a\to c\to b\) for the fact
\((a,b,1)\), and \(b\to d_1\to d_2\to b\) for the fact \((b,b,2)\).  The looped
vertices are \(a,b\).  \(\pi_1(a,b)\) holds through \(c\), and \(\pi_2(b,b)\)
through \(d_1,d_2\).  \(\pi_1(b,a)\) fails: the only loop-free vertex reached
from \(b\) is \(d_1\), whose only successor \(d_2\) is loop-free.

\emph{The refined formula.}  Let \(M\) have \(\delta(q_0,a)=(q,a,R)\) for every
tape symbol \(a\), with a target state \(q\) fixed in three variants.
\begin{itemize}
\item \(q=q_{\mathrm{acc}}\): the computation accepts after \(k=1\) step.  The
  model of Theorem~\ref{thm:p02-w02-p02-036-refined}(2) has domain
  \(\{-2,\ldots,2\}\) with \(K\) the usual order, \(S\) the successor pairs
  \((x,x+1)\), \(Z=\{0\}\), \(I_{q_0}=\{0\}\), \(I_{q_{\mathrm{acc}}}=\{1,2\}\),
  \(H=\{(0,0),(1,1),(2,1)\}\), and \(T_{\sqcup}\) holding at every pair
  \((t,x)\) with \(t\in\{0,1,2\}\).  The moment \(2\) repeats moment \(1\).  The
  only active moment is \(0\), whose successor is \(1\), and the visited cells
  \(0\) and \(1\) have both neighbours in the domain.  It has \(2k+3=5\)
  elements.
\item \(q=q_{\mathrm{rej}}\): in any model, the element \(z\) with \(Z(z)\)
  carries \(I_{q_0}\), so it is active and \(\prec_{K,M}\) gives a successor
  moment \(t'\); the visited cell \(z\) gets a successor \(x'\).  The step
  clause for \((q_0,\sqcup)\) then gives \(I_{q_{\mathrm{rej}}}(t')\), against
  \(\mathrm{NONSTOP}\).  So \(\rho_{M,\prec}\) is contradictory.
\item \(q=q_0\): the head moves right forever.  The integers, with every
  moment \(t\geq0\) in state \(q_0\) and head at cell \(t\), form a model, and
  by Lemma~\ref{thm:p02-w02-p02-036-simulation} every model contains the
  pairwise distinct moments \(\tau_0,\tau_1,\ldots\).  So \(\rho_{M,\prec}\) is
  an infinity axiom.
\end{itemize}
\end{hbexample}

\begin{hbwarning}{def:p02-w02-p02-036-limits}{What the theorems do not give}
\begin{enumerate}
\item The binary relation symbol is used in an essential way: the path
  translation needs it to code every relation of the run signature and of
  \(\Lang_M^{\prec}\).  The hypothesis cannot be deleted.  If every relation
  symbol of \(\Lang\) is unary, every \(\Lang\)-sentence lies in the class
  \([\mathit{all},(\omega),(1)]_{=}\) of sentences with arbitrary prefix,
  unary predicates, at most one unary function symbol and equality, in the
  prefix-vocabulary notation of
  \cite[Sect.~1.3, pp.~8--9]{BorgerGradelGurevich1997}.  Finite
  satisfiability for that class is decidable
  \cite[Thm.~6.0.2(6), p.~240]{BorgerGradelGurevich1997}.
\item Finite validity is co-computably enumerable: a search for a finite
  countermodel halts exactly on the sentences that are \emph{not} finitely
  valid.  What fails is the converse direction.  No effective calculus is
  sound and complete for truth in all finite \(\Lang\)-structures, since its
  theorems would form a computably enumerable set equal to the finite
  validities.
\item The results concern one fixed signature, or all formulas, and the full
  set of sentences.  They do not say that any particular sentence is hard to
  classify, and they give no bound for the size of a finite model when one
  exists.  By Corollary~\ref{thm:p02-w02-p02-036-trakhtenbrot} no computable
  bound exists, since such a bound would decide \(\mathrm{Fin\text{-}sat}(X)\)
  by a finite search through the structures of at most that size.
\item The finite model property is a sufficient condition for decidability
  of satisfiability, not a necessary one.  The one-element class
  \(\{\prec_K\}\) is decidable, its satisfiability problem is decidable
  because its only member is satisfiable, and it lacks the finite
  model property by Lemma~\ref{thm:p02-w02-p02-036-infinity-axiom}.
\item Unrestricted compactness does not transfer to finite structures
  (\cref{def:p03-finite-compactness-failure}), and the first-order zero--one
  law (\cref{thm:p03-fol-zero-one}) supplies asymptotic probabilities, not a
  decision procedure for finite validity.
\end{enumerate}
\end{hbwarning}

\emph{Relationships.}  The proofs use the machine model of
\cref{sec:turing-machines-01}, the program codes of
\cref{sec:universal-reductions-01}, the undecidability of \(K\)
(\cref{thm:halting-undecidable}), and the complement criterion of
\cref{sec:turing-machines-03}.  Unrestricted validity contrasts with them
through \cref{thm:fol-validity-enumerable}.  The finite-model-theory survey
uses the theorems in \cref{sec:p03-finite-unrestricted}.  Two reductions
appear.  The run sentence of \cref{sec:p02-w02-p02-036-run-sentence} controls
finite models only and yields Trakhtenbrot's theorem in its undecidability
form.  The refined formula of \cref{sec:p02-w02-p02-036-refined}, the
relational version of the source's \(\rho_{M,\prec}\), controls all models and
yields the inseparable form.  The path translation and the renaming carry both
to any finite relational signature with a binary symbol, and they are
conservative reductions in the sense of
Definition~\ref{def:p02-w02-p02-036-conservative}.

\emph{Sources.}  The notation and the three classes are those of
\cite[Sect.~2.1.2, Def.~2.1.27]{BorgerGradelGurevich1997}; the finite model
property argument is on p.~33, the infinity axiom \(\prec_K\) and
Exercise~2.1.26 on p.~33, Definition~2.1.28 and Exercise~2.1.29 on p.~33,
Theorem~2.1.30 with its proof on pp.~34--35, and Definition~2.1.35 and
Exercise~2.1.36 on p.~37; the Church--Turing reduction formula \(\rho_M\) and
the step formulas (2.1)--(2.3) are in Sect.~2.1.1, pp.~20--22.  The result is
due to Trakhtenbrot \cite{Trakhtenbrot1950}.  The proofs here are written for
this handbook's machine model.  Where the source sketches a step, in
particular the verification of the refined reduction property, the annex gives
the full argument for the relational formula of
Definition~\ref{def:p02-w02-p02-036-refined}, whose differences from the
source's formula are stated after that definition.
